\chapter{四维超对称性与Wess--Zumino模型}

\lectureoverview{上一讲的一维超对称代数只涉及哈密顿量。洛伦兹不变性现在迫使能量与空间动量重新结合，而超荷必须成为二分量旋量。我们通过首先回到无质量Weyl费米子及其哈密顿量\(H=\boldsymbol{\sigma}\cdot\mathbf p\)来得到该代数。随后我们将超荷表示为四维超空间上的微分算符，用约束\(\bar D_{\dot\alpha}\Phi=0\)分离出手征超场，并利用超空间积分构造一个具体的场论。依次选择\(\Phi^\dagger\Phi\)、\(\Phi^2\)和\(\Phi^3\)会产生传播、玻色子--费米子质量相等、Yukawa相互作用以及一个相关的标量相互作用。辅助场\(F\)使这些关系变得透明，而结尾的圈图则表明超对称为何会软化紫外敏感性。}

\section{为什么哈密顿量不能单独存在}

在玩具模型中，空间被抑制，基本关系为
\begin{equation}
\{Q,Q^\dagger\}=2H.
\label{eq:l8-time-only-algebra}
\end{equation}
当时间坐标是唯一的坐标时，这种形式是足够的。然而在相对论性理论中，\(H=P_0\)是四维动量\(P_\mu\)的一个分量。洛伦兹助推会将它和\(P_i\)混合。一个右侧单独挑出\(P_0\)的代数不可能协变地变换。

因此，修补必须同时做两件事。它必须用全部四个分量\(P_\mu\)替换\(H\)，并且必须在左侧提供能够与右侧洛伦兹矢量相匹配的指标。无质量Dirac方程提供了所需的线索：相对论性费米子可以由二分量旋量描述，而矩阵\(\sigma^\mu\)将矢量指标转换为自旋指标对。

\section{无质量费米子、螺旋度与手征性}

从有质量粒子与无质量粒子之间的物理区别开始。取一个自旋为\(\tfrac12\)的有质量粒子沿其自旋方向运动。由于其速率小于光速，观察者可以追上它。在新参考系中，动量反向而自旋不变，因此螺旋度反向。因此，洛伦兹不变的有质量理论必须同时包含两种螺旋度。

没有观察者能追上无质量粒子。固有洛伦兹变换不能使其动量相对于自旋反向，因此一个不可约的无质量表示可以只携带一种螺旋度。空间旋转可以使飞行方向反转，但它不会把右手变成左手。

\begin{classroomqa}
\audiencequestion{追及论证讨论的究竟是螺旋度而不是手征性吗？\lecturetimestamp{00:05:42}}
\lecturerresponse{是的。螺旋度是自旋沿动量方向的投影，而手征性是手征矩阵的本征值，通常记为\(\gamma^5\)。对于无质量模式，这两个标签绑定在一起，这就是为什么在这个例子中这两种说法听起来可以互换。对于有质量粒子，它们是不同的。}
\end{classroomqa}

态计数需要一个限定。\editorialnote{单个无质量Weyl场具有一种粒子螺旋度和相反的反粒子螺旋度。因此，“一半的态数”指的是从四分量Dirac场约化为一个二分量Weyl场；这并不意味着反粒子态消失了。}

将本开头论证中使用的Weyl场写为
\begin{equation}
\chi=
\begin{pmatrix}
\chi_1\\
\chi_2
\end{pmatrix},
\end{equation}
本讲用\(\psi_\alpha\)表示这个场；我们在这里称它为\(\chi\)，以免与下面通常置于手征超场中的共轭Weyl表示混淆。
\begin{equation}
H\chi=\boldsymbol{\sigma}\cdot\mathbf p\,\chi.
\label{eq:l8-weyl-hamiltonian}
\end{equation}
对于动量\(\mathbf p\)，矩阵\(\boldsymbol{\sigma}\cdot\mathbf p\)具有本征值\(\pm|\mathbf p|\)。正频支提供粒子模式；负频支通过反粒子构造重新解释，而不是作为无下界的谱保留下来。

\begin{figure}[H]
\centering
\includegraphics[width=0.70\textwidth]{lecture_08_figure_02.png}
\lectureframe{00:07:44.271}
\caption{手征二分量旋量与哈密顿量\(H=\boldsymbol{\sigma}\cdot\mathbf p\)。最后的动量符号被部分遮挡，但口述的方程确定了它。}
\end{figure}

手性标签也依赖于约定。\editorialnote{黑板将一种手性直接与\(\boldsymbol{\sigma}\cdot\mathbf p\)的正本征值联系起来。哪种手征性获得该符号取决于\(\sigma^\mu\)、\(\bar\sigma^\mu\)和Weyl场的定义。我们保留本征值陈述，但不给它附加与约定无关的左/右标签。}

将 \(H\) 替换为 \(i\partial_t\)，并将 \(\mathbf p\) 替换为 \(-i\nabla\)，得到
\begin{equation}
i\partial_t\chi=-i\boldsymbol{\sigma}\cdot\nabla\chi.
\end{equation}
引入
\begin{equation}
\begin{aligned}
\sigma^\mu&=(\mathbf 1,\boldsymbol{\sigma}),\\
\bar\sigma^\mu&=(\mathbf 1,-\boldsymbol{\sigma}),\\
\eta_{\mu\nu}&=\operatorname{diag}(1,-1,-1,-1).
\end{aligned}
\label{eq:l8-sigma-convention}
\end{equation}
于是刚刚得到的方程成为
\begin{equation}
i\sigma^\mu\partial_\mu\chi=0.
\label{eq:l8-weyl-equation}
\end{equation}
共轭 Weyl 表示使用 \(\bar\sigma^\mu\)。正是这个共轭表示——下文称为 \(\psi_\alpha\)——其自由方程为 \(i\bar\sigma^\mu\partial_\mu\psi=0\)。证明协变性需要 Lorentz 群的旋量表示，但这里重要的事实是结构性的：四个矩阵 \(\sigma^\mu\) 将一个时空指标与两个旋量指标连接起来。

\section{四维超代数}

奇坐标现在必须像 Weyl 旋量那样变换：
\begin{equation}
\theta^\alpha=(\theta^1,\theta^2),
\qquad
\bar\theta^{\dot\alpha}=(\bar\theta^{\dot1},\bar\theta^{\dot2}).
\end{equation}
横线表示复共轭，点用于区分共轭 Lorentz 表示。共有两个复的、即四个实的 Grassmann 分量。因此，一般的偶超场
\begin{equation}
\mathcal S=\mathcal S(x^\mu,\theta^\alpha,\bar\theta^{\dot\alpha})
\end{equation}
具有有限的 Grassmann 展开。其最高单项式为
\begin{equation}
\theta^2\bar\theta^2
\equiv
(\theta^1\theta^2)(\bar\theta^{\dot1}\bar\theta^{\dot2}),
\end{equation}
即全部四个独立奇坐标之积。该单项式的系数通常称为 \(D\) 分量。

超荷也带有旋量指标。在四维 \(N=1\) 超对称中，其代数为
\begin{align}
\{Q_\alpha,\bar Q_{\dot\beta}\}
&=2\sigma^\mu_{\alpha\dot\beta}P_\mu,
\label{eq:l8-superalgebra-main}\\
\{Q_\alpha,Q_\beta\}&=0,
&
\{\bar Q_{\dot\alpha},\bar Q_{\dot\beta}\}&=0,
\label{eq:l8-superalgebra-nilpotent}\\
[P_\mu,Q_\alpha]&=0,
&
[P_\mu,\bar Q_{\dot\alpha}]&=0,
&
[P_\mu,P_\nu]&=0.
\label{eq:l8-superalgebra-commutators}
\end{align}
因子 2 是约定俗成的。矩阵
\(\sigma^\mu_{\alpha\dot\beta}\)
缩并 \(P_\mu\) 上的矢量指标，并恰好留下出现在左侧的两个旋量指标。

\begin{classroomqa}
\audiencequestion{\(\alpha\) 和 \(\beta\) 是 sigma 矩阵上的矩阵指标吗？\lecturetimestamp{00:22:01}}
\lecturerresponse{是的。更确切地说，其中一个是无点 Weyl 指标，另一个是有点共轭指标；\(\mu\) 是时空矢量指标。sigma 矩阵是这些描述之间的桥梁。}
\end{classroomqa}

整理后的记号使这一桥梁变得明确。\editorialnote{讲演有意省略了点指标，并在黑板上讨论 \(\alpha\beta\) 与 \(\beta\alpha\) 的次序。方程 ~\eqref{eq:l8-superalgebra-main} 和 \eqref{eq:l8-superalgebra-commutators} 恢复了标准的有点/无点记号，而不改变所论证的内容。}

最后一行说明超荷是守恒的，并且与平移相容。动量分量对易，因为它们的导数表示对易。在下面使用的微分约定中，
\begin{equation}
P_\mu=i\partial_\mu,
\end{equation}
因此式 ~\eqref{eq:l8-superalgebra-main} 也可以写成
\begin{equation}
\{Q_\alpha,\bar Q_{\dot\beta}\}
=2i\sigma^\mu_{\alpha\dot\beta}\partial_\mu.
\end{equation}

\section{作为超空间导数的超荷}

动量平移普通坐标；超荷平移奇坐标，同时也移动时空。固定左 Grassmann 导数并定义
\begin{align}
Q_\alpha
&=\partial_\alpha
-i\sigma^\mu_{\alpha\dot\alpha}
\bar\theta^{\dot\alpha}\partial_\mu,
\label{eq:l8-Q}\\
\bar Q_{\dot\alpha}
&=-\bar\partial_{\dot\alpha}
+i\theta^\alpha\sigma^\mu_{\alpha\dot\alpha}\partial_\mu.
\label{eq:l8-Qbar}
\end{align}
用 \(\epsilon^\alpha Q_\alpha+
\bar\epsilon_{\dot\alpha}\bar Q^{\dot\alpha}\) 作用，给出超场的无穷小超对称变分。其分量就是各独立 Grassmann 单项式的系数，而微分算子混合这些系数。

\begin{classroomqa}
\audiencequestion{在简写 \(\bar\theta\sigma\partial\) 中指的是哪个 \(x\) 导数？\lecturetimestamp{00:25:40}}
\lecturerresponse{全部四个。\(\bar\theta^{\dot\alpha}\sigma^\mu_{\alpha\dot\alpha}\partial_\mu\) 中重复的 \(\mu\) 被求和，因此 \(\sigma^\mu\) 的每个分量都与相应的导数 \(\partial/\partial x^\mu\) 配对。}
\end{classroomqa}

\begin{classroomqa}
\audiencequestion{共轭荷应当包含 \(\theta\)，而不是另一个 \(\bar\theta\) 吗？\lecturetimestamp{00:26:16}}
\lecturerresponse{是的。\(Q\) 对 \(\theta\) 求导，并包含 \(\bar\theta\)；\(\bar Q\) 对 \(\bar\theta\) 求导，并包含 \(\theta\)。它们是共轭算子。}
\end{classroomqa}

\begin{classroomqa}
\audiencequestion{\(\sigma^\mu\) 是对称的吗？\lecturetimestamp{00:27:04}}
\lecturerresponse{Pauli 矩阵是 Hermite 的。\((\mathbf1,\boldsymbol{\sigma})\) 中四个矩阵里有三个是对称的，而 \(\sigma^2\) 在普通转置下是反对称的。Hermite 性，连同有点与无点指标的放置，才是相关的共轭性质。}
\end{classroomqa}

\begin{classroomqa}
\audiencequestion{共轭方程中的 sigma 矩阵是否以相反次序携带指标？\lecturetimestamp{00:27:13}}
\lecturerresponse{共轭交换有点与无点表示，因此稳妥的记号是 \(\sigma^\mu_{\alpha\dot\alpha}\) 及其共轭 \(\bar\sigma^{\mu\,\dot\alpha\alpha}\)，而不是把这两个标签当作可互换的矩阵槽位。}
\end{classroomqa}

式 ~\eqref{eq:l8-Q} 和 \eqref{eq:l8-Qbar} 中的公式与黑板上使用的简写在外观符号上不同。这些符号与 \(P_\mu\) 是 \(+i\partial_\mu\) 还是 \(-i\partial_\mu\)、导数是左作用还是右作用，以及共轭奇变量如何排序有关。重要的是，一种约定要同时满足整个代数。

\section{一个在超对称下仍然成立的约束}

一般的超场包含许多分量。我们希望在不破坏对称性的前提下寻找一个更小的类。本讲先给出一个普通的类比。对于矢量\(\mathbf V\)，约束
\begin{equation}
V_x=0
\end{equation}
不是旋转不变的：旋转之后，\(V_y\)会对新的\(x\)分量有贡献。相比之下，
\begin{equation}
\mathbf V^2=1
\end{equation}
是不变的，并且去掉一个独立自由度。一个有用的约束必须在对称性下把受约束的场映射到受约束的场。

假设一个线性奇算子\(\bar D_{\dot\alpha}\)定义
\begin{equation}
\bar D_{\dot\alpha}\Phi=0.
\label{eq:l8-chiral-constraint}
\end{equation}
在一次超对称变分\(\delta\Phi=Q\Phi\)之后，如果\(\bar D\)与超荷反交换，同一条件仍将成立：
\begin{equation}
\bar D(Q\Phi)=-Q(\bar D\Phi)=0.
\end{equation}
这是促使我们进行构造的抽象一致性检验，而不是一个不加解释就出现的定义。

\begin{classroomqa}
\audiencequestion{为什么在最后一个表达式中\(Q\)和\(\bar D\)可以交换？\lecturetimestamp{00:36:56}}
\lecturerresponse{它们是选来反交换的奇算子。因此交换它们的顺序会引入一个负号；一旦\(\bar D\)作用到\(\Phi\)上，原来的约束就使结果为零。}
\end{classroomqa}

所需的协变导数通过相对于超荷反转时空导数符号而得到：
\begin{align}
D_\alpha
&=\partial_\alpha
+i\sigma^\mu_{\alpha\dot\alpha}
\bar\theta^{\dot\alpha}\partial_\mu,
\label{eq:l8-D}\\
\bar D_{\dot\alpha}
&=-\bar\partial_{\dot\alpha}
-i\theta^\alpha\sigma^\mu_{\alpha\dot\alpha}\partial_\mu.
\label{eq:l8-Dbar}
\end{align}
它们满足
\begin{align}
\{D_\alpha,\bar D_{\dot\beta}\}
&=-2i\sigma^\mu_{\alpha\dot\beta}\partial_\mu,\\
\{D_\alpha,Q_\beta\}
&=\{D_\alpha,\bar Q_{\dot\beta}\}=0,\\
\{\bar D_{\dot\alpha},Q_\beta\}
&=\{\bar D_{\dot\alpha},\bar Q_{\dot\beta}\}=0.
\end{align}

\begin{classroomqa}
\audiencequestion{从\(\bar Q\)到\(\bar D\)的唯一变化是时空导数项的符号吗？\lecturetimestamp{00:37:58}}
\lecturerresponse{正是如此。这一符号反转使得当\(D\)与任一超荷反交换时，交叉项相互抵消。}
\end{classroomqa}

我们只施加式~\eqref{eq:l8-chiral-constraint}。同时施加\(D_\alpha\Phi=0\)和\(\bar D_{\dot\alpha}\Phi=0\)会通过它们的反交换子迫使\(\partial_\mu\Phi=0\)，至多留下一个与时空无关的剩余部分。
\editorialnote{黑板上说，同时施加两个约束会迫使\(\Phi=0\)。代数严格地迫使时空无关性；额外的边界条件或分量条件随后可能进一步约化该场。}

\section{求解手征约束}

定义平移后的坐标
\begin{equation}
y^\mu=x^\mu+i\theta\sigma^\mu\bar\theta.
\label{eq:l8-chiral-coordinate}
\end{equation}
在上述左导数约定下，函数
\begin{equation}
\Phi=\Phi(y,\theta)
\label{eq:l8-chiral-solution}
\end{equation}
满足\(\bar D_{\dot\alpha}\Phi=0\)。链式法则使抵消变得明确。对\(y^\mu\)求导会给出同一个\(i\theta\sigma^\mu\partial_\mu\Phi\)项，而它在\(\bar D_{\dot\alpha}\)中以相反符号出现。

因此，手征超场没有独立的\(\bar\theta\)依赖。它仍然可以包含\(\bar\theta\)，但只能通过组合\(y^\mu\)包含。其有限展开为
\begin{equation}
\Phi(y,\theta)
=\phi(y)+\sqrt2\,\theta\psi(y)+\theta^2F(y).
\label{eq:l8-chiral-expansion-y}
\end{equation}
三种分量类型分别是一个复标量\(\phi\)、一个Weyl费米子\(\psi_\alpha\)和一个复标量\(F\)。最后一个是辅助场，而不是第二个传播玻色子。

\begin{classroomqa}
\audiencequestion{\(x\)是否已被一个与它正交的新坐标\(y\)取代？\lecturetimestamp{00:49:08}}
\lecturerresponse{不是。\(y^\mu\)是式~\eqref{eq:l8-chiral-coordinate}中那个特定组合的简写；它并不是一个额外的时空方向。在\(y\)中工作使手征性变得简单，而在构造普通场论时我们再翻译回\(x\)。}
\end{classroomqa}

在固定\(x\)下展开式~\eqref{eq:l8-chiral-expansion-y}给出
\begin{align}
\Phi(x,\theta,\bar\theta)
={}&\phi
+\sqrt2\,\theta\psi
+\theta^2F
+i\theta\sigma^\mu\bar\theta\,\partial_\mu\phi
\notag\\
&-\frac{i}{\sqrt2}\theta^2
\partial_\mu\psi\,\sigma^\mu\bar\theta
+\frac14\theta^2\bar\theta^2\Box\phi.
\label{eq:l8-chiral-expansion-x}
\end{align}
右端所有分量场都在\(x\)处取值。共轭反手征场依赖于
\(\bar y^\mu=x^\mu-i\theta\sigma^\mu\bar\theta\)和\(\bar\theta\)。手征超场不可能自洽地为实，因为它的共轭满足相反的手征约束。

\section{超空间积分}

全超空间作用量具有形式
\begin{equation}
S_D=\int d^4x\,d^2\theta\,d^2\bar\theta\;K,
\label{eq:l8-D-term-action}
\end{equation}
其中\(K\)是实超场。Berezin积分选出其
\(\theta^2\bar\theta^2\)系数。在奇积分之后，只剩下普通的时空拉格朗日量。

还有第二种不变构造。如果\(W(\Phi)\)是手征的，那么
\begin{equation}
S_F=\int d^4x\,d^2\theta\;W(\Phi)
+\int d^4x\,d^2\bar\theta\;\overline{W}(\Phi^\dagger)
\label{eq:l8-F-term-action}
\end{equation}
是超对称的。第一个积分选出\(\theta^2\)分量，第二个积分选出其共轭。

\begin{classroomqa}
\audiencequestion{当表达式没有\(\bar\theta\)依赖时，为什么只对\(d^2\theta\)积分是合法的？\lecturetimestamp{01:16:30}}
\lecturerresponse{手征性在超对称下保持不变，而手征超场的最高\(\theta\)分量只按时空全导数变化。因此它的时空积分是不变的，正如完整的\(D\)分量积分是不变的一样。}
\end{classroomqa}

\section{自由手征多重态}

取这个看似没有结构的表达式
\begin{equation}
S_{\mathrm{free}}
=\int d^4x\,d^2\theta\,d^2\bar\theta\;
\Phi^\dagger\Phi.
\label{eq:l8-free-superaction}
\end{equation}
两个因子依赖于相反平移的坐标：
\begin{equation}
\Phi=\Phi(x+i\theta\sigma\bar\theta,\theta),
\qquad
\Phi^\dagger
=\Phi^\dagger(x-i\theta\sigma\bar\theta,\bar\theta).
\end{equation}
在时空积分内，我们可以平移哑变量\(x\)。一次平移可以去掉一个因子中的Grassmann位移，但随后会使另一个因子中的相对位移加倍。对该幂零位移作Taylor展开会在有限多项后终止。

导数正是从这里来的。超空间表达式
\(\Phi^\dagger\Phi\)
不包含任何显式的时空导数，然而平移后的自变量却生成了各分量场的一阶和二阶导数。将展开式相乘并只保留最顶层的分量，按式~\eqref{eq:l8-sigma-convention}的约定，得到
\begin{equation}
\mathcal L_{\mathrm{free}}
=\partial_\mu\phi^\ast\partial^\mu\phi
+i\bar\psi\bar\sigma^\mu\partial_\mu\psi
+F^\ast F.
\label{eq:l8-free-component-lagrangian}
\end{equation}
在分部积分的意义下，标量部分也可以示意性地写成\(-\phi^\ast\Box\phi\)，这正是板书计算中所采用的形式。

\begin{classroomqa}
\audiencequestion{超场泰勒展开中的第二行是从哪里来的？\lecturetimestamp{01:03:00}}
\lecturerresponse{它来自分量展开\(\Phi(y,\theta)=\phi(y)+\sqrt2\theta\psi(y)+\theta^2F(y)\)。随后每个分量都要做泰勒展开，因为\(y\)与\(x\)相差一个幂零的格拉斯曼双线性项。}
\end{classroomqa}

\begin{classroomqa}
\audiencequestion{最后那个分量应该带有\(\theta^2\)吗？\lecturetimestamp{01:04:54}}
\lecturerresponse{是的。辅助分量是\(F\theta^2\)。由于它已经含有两个独立的\(\theta\)，任何再携带一个\(\theta\)的进一步泰勒平移都会消失。}
\end{classroomqa}

\begin{figure}[H]
\centering
\includegraphics[width=0.82\textwidth]{lecture_08_figure_03.png}
\lectureframe{01:08:57.900}
\caption{平移坐标展开被归入普通拉格朗日量的标量项、费米子项和辅助项。}
\end{figure}

式~\eqref{eq:l8-free-component-lagrangian}中的前两项给出无质量的克莱因--戈尔登方程和韦尔方程。最后一项不含导数：
\begin{equation}
\frac{\partial\mathcal L}{\partial F}=F^\ast=0,
\qquad
\frac{\partial\mathcal L}{\partial F^\ast}=F=0.
\label{eq:l8-free-F-equation}
\end{equation}
它不引入任何初始数据，也不对应任何在壳粒子。

\begin{figure}[H]
\centering
\includegraphics[width=0.82\textwidth]{lecture_08_figure_04.png}
\lectureframe{01:11:32.360}
\caption{分量拉格朗日量末尾那个不含导数的辅助项\(F^\ast F\)。}
\end{figure}

传播子的语言说的是同一件事。\(F\)的逆二次算符与动量无关。它的傅里叶变换是局域的：
\begin{equation}
\langle F(x)F^\ast(y)\rangle
\propto \delta^{(4)}(x-y).
\label{eq:l8-F-contact}
\end{equation}
一条\(F\)线可以组织一个图，但它的两个端点重合。它是一种接触收缩，而不是物理粒子的轨迹。

\section{二次超势与相等质量}

传播来自\(D\)项。质量来自手征积分。取常规归一化
\begin{equation}
W_2(\Phi)=\frac12m\Phi^2.
\end{equation}
\(\theta^2\)分量为
\begin{equation}
\left.W_2(\Phi)\right|_{\theta^2}
=m\phi F-\frac12m\psi\psi.
\end{equation}
与共轭项一起，分量拉格朗日量变为
\begin{align}
\mathcal L
={}&
\partial_\mu\phi^\ast\partial^\mu\phi
+i\bar\psi\bar\sigma^\mu\partial_\mu\psi
+F^\ast F
\notag\\
&+m\phi F+m^\ast\phi^\ast F^\ast
-\frac12m\psi\psi
-\frac12m^\ast\bar\psi\bar\psi.
\label{eq:l8-mass-offshell}
\end{align}
费米子双线性项是韦尔场的马约拉纳质量。\(F\)方程为代数方程：
\begin{equation}
F^\ast=-m\phi,
\qquad
F=-m^\ast\phi^\ast.
\label{eq:l8-mass-F-equations}
\end{equation}
代入后得到
\begin{equation}
\begin{aligned}
\mathcal L_{\mathrm{on\ shell}}={}&
\partial_\mu\phi^\ast\partial^\mu\phi
+i\bar\psi\bar\sigma^\mu\partial_\mu\psi\\
&-|m|^2|\phi|^2\\
&-\frac12\left(m\psi\psi+m^\ast\bar\psi\bar\psi\right).
\end{aligned}
\label{eq:l8-mass-onshell}
\end{equation}
因此标量质量和费米子质量都是\(|m|\)。用板书上的图示语言来说，一个标量通过一次\(m\phi F\)插入局域地变成\(F\)，并立即通过其共轭变回来。两次插入给出\(|m|^2|\phi|^2\)，正是通过求解式~\eqref{eq:l8-mass-F-equations}所得到的标量质量项。

现在可以精确地陈述归一化了。\editorialnote{讲座在相乘超场时有意省略了数值系数。这里使用的因子\(1/2\)、\(1/3\)和\(\sqrt2\)是常规的韦斯--祖米诺归一化；它们使分量质量和耦合精确，同时保留板书推导的每一个定性步骤。}

\section{三次超势与相互作用}

现在加上
\begin{equation}
W(\Phi)=\frac12m\Phi^2+\frac13g\Phi^3.
\label{eq:l8-superpotential}
\end{equation}
对于一般的全纯\(W\)，\(F\)分量为
\begin{equation}
\left.W(\Phi)\right|_{\theta^2}
=F W'(\phi)-\frac12W''(\phi)\psi\psi.
\label{eq:l8-W-components}
\end{equation}
这里
\begin{equation}
W'(\phi)=m\phi+g\phi^2,
\qquad
W''(\phi)=m+2g\phi.
\end{equation}
在消去\(F\)之前，三次项含有
\begin{equation}
g\phi^2F-g\phi\psi\psi+\text{h.c.}
\label{eq:l8-cubic-offshell}
\end{equation}
第二项是一个汤川顶点，将一个标量与两条费米子线连接起来。第一项将两个标量与一条辅助线连接起来。

\begin{figure}[H]
\centering
\begin{tikzpicture}[scale=0.95,>=Latex]
  \coordinate (vone) at (0,0);
  \draw[->,thick] (-2,1.1) -- (vone);
  \draw[->,thick] (vone) -- (-2,-1.1);
  \draw[thick,dashed] (vone) -- (2,0);
  \node[left] at (-2,1.1) {\(\psi\)};
  \node[left] at (-2,-1.1) {\(\psi\)};
  \node[right] at (2,0) {\(\phi\)};
  \node[below=8pt of vone] {\(-g\phi\psi\psi\)};
\end{tikzpicture}
\caption{汤川相互作用的教学重绘：两条费米子腿和一条标量腿相交，耦合正比于\(g\)。}
\end{figure}

完整的辅助场方程为
\begin{equation}
F^\ast=-W'(\phi)=-(m\phi+g\phi^2).
\label{eq:l8-general-F-equation}
\end{equation}
消去它后得到标量势
\begin{equation}
V(\phi)=|W'(\phi)|^2
=|m\phi+g\phi^2|^2.
\label{eq:l8-scalar-potential}
\end{equation}
除质量项外，它还含有正比于\(m^\ast g\)的三次标量项以及一个四次项
\begin{equation}
V(\phi)\supset |g|^2|\phi|^4.
\label{eq:l8-quartic}
\end{equation}
在辅助线图像中，两个\(\phi^2F\)顶点通过局域的\(F\)收缩相连。由于式~\eqref{eq:l8-F-contact}迫使它们的位置重合，结果正是一个四标量接触顶点。

\begin{figure}[H]
\centering
\begin{tikzpicture}[scale=0.95]
  \coordinate (vfour) at (0,0);
  \foreach \angle in {35,145,215,325}
    \draw[thick,dashed] (vfour) -- (\angle:1.9);
  \fill (vfour) circle (1.5pt);
  \node[below=8pt of vfour] {\(-|g|^2|\phi|^4\)};
\end{tikzpicture}
\caption{消去辅助场后，两个局域\(F\)顶点变成一个四标量接触相互作用。}
\end{figure}

超对称性所做的不仅仅是将一个玻色子和一个费米子放进同一个理论中。它迫使它们的质量一致，将汤川耦合与\(g\)联系起来，并把四次标量耦合固定为\(|g|^2\)。这些关系可以通过质量、散射振幅和衰变来检验；它们不是可有可无的数值巧合。

\section{成对的圈效应}

当相互作用修正粒子质量时，这些关系变得最为有用。标量自能通过四顶点获得一个玻色子圈，并通过两个汤川顶点获得一个费米子圈。每个都带有同一耦合常数的两次幂。闭合费米子圈具有费米统计所要求的相对负号。

\begin{figure}[H]
\centering
\begin{tikzpicture}[>=Latex,scale=0.9]
  \begin{scope}[xshift=-3.3cm]
    \draw[thick,dashed] (-2,0)--(0,0)--(2,0);
    \draw[thick,dashed] (0,0) .. controls (-0.9,1.7) and (0.9,1.7) .. (0,0);
    \fill (0,0) circle (1.4pt);
    \node at (0,-0.55) {bosonic loop};
  \end{scope}
  \begin{scope}[xshift=3.3cm]
    \draw[thick,dashed] (-2,0)--(-0.8,0);
    \draw[thick,dashed] (0.8,0)--(2,0);
    \draw[->,thick] (-0.8,0) arc[start angle=180,end angle=0,x radius=0.8,y radius=0.8];
    \draw[->,thick] (0.8,0) arc[start angle=0,end angle=-180,x radius=0.8,y radius=0.8];
    \fill (-0.8,0) circle (1.4pt);
    \fill (0.8,0) circle (1.4pt);
    \node at (0,-1.35) {fermionic loop};
  \end{scope}
\end{tikzpicture}
\caption{讲义中成对的标量自能机制。超对称性固定了相对耦合，而闭合费米子圈提供了相反的符号。}
\end{figure}

在一圈阶，领头紫外敏感部分在完整的玻色子与费米子贡献之间相消。同一对称性通过沃德恒等式和非重整化性质组织更高阶项。该陈述涉及与对称性相关的求和，而不是每个图逐一的字面相消；波函数重整化、有限效应以及其他保持对称性的修正都可以保留。这正是用超空间构建理论的实用威力所在：质量的相等以及耦合之间的关系在任何单个费曼积分被计算之前就已被强制成立。\editorialnote{黑板上更宽泛地陈述了这种相消；正文给出了标准的限定条件。}

\section{接下来是什么}

这里构建的模型是最简单的相互作用手征理论，通常称为Wess--Zumino模型。它是一个例子，而不是最一般的超对称场论，但其模式仍然成立：用超多重态表示场，写出不变的超空间积分，并让格拉斯曼展开强制给出各分量之间的关系。

精确超对称性在观测到的粒子谱中并不可见，因此下一个问题是如何在保留足够多的紫外控制以使其有用的同时破缺该对称性。另一个所承诺的方向是大统一：如何
\begin{equation}
SU(3)\times SU(2)\times U(1)
\end{equation}
可能嵌入到诸如\(SU(5)\)这样的更大群中，以及超对称粒子内容如何改变那个统一故事。这些就是剩余讲义所转向的问题。
